\section{Threats to Validity}

\para{Internal validity.}
Our code detection heuristic (code ratio $>$ 20\%, length $>$ 2{,}000 characters) is a simple rule that may produce false positives (rendering non-code outputs as images) or false negatives (missing code outputs). In our 20-instance pilot, only 3 of 92 rendered outputs were borderline cases (Python script outputs), and manual inspection confirmed they contained code-like content. A more sophisticated classifier could improve precision, but would introduce its own complexity and potential confounds.

The image rendering pipeline introduces design choices---font size, page width, line wrapping---that were fixed via the pilot study. A different specification might yield different results. We mitigate this by selecting parameters that achieve near-perfect fidelity ($\le$1.5\% CER, 100\% identifier recall) and documenting the selection process transparently.

GPT-5-mini does not support \texttt{temperature=0}, so all outputs involve sampling noise. The 10/10 symmetric disagreement between conditions may partly reflect this randomness rather than genuine modality differences. Multiple runs per instance would be needed to disentangle these effects.

\para{External validity.}
We evaluate a single model (GPT-5-mini) on a single benchmark (100-instance subset of SWE-bench Verified). Our findings may not generalize to other vision-capable LLMs with different vision encoders or to other coding tasks beyond patch generation. SWE-bench Verified is predominantly Python; optical encoding may behave differently for other languages with different visual density or indentation conventions.

\para{Construct validity.}
The SWE-bench pass rate is binary: a patch either passes all required tests or fails. This does not capture partial correctness or near-miss patches. Our cost comparison depends on current OpenAI API pricing, which may change. The code detection heuristic's thresholds (2{,}000 chars, 20\% code ratio) were chosen empirically and not optimized; different thresholds could shift the balance of what gets rendered.
